\section{Week 12}

\subsection{Opening Remarks}

Recall that the powerful idea that we encountered from last week is the mapping \( {\pi}_{H} : G \to {S}_{G/H} \) especially \( \text{ker}({\pi}_{H}) = \bigcap_{  g \in G  }^{  }  g H g^{-1} \leq H  \). If \( G  \) is a group and \( H  \) is a subgroup with \(  | G : H  |  = p  \) where \( p  \) is the smallest prime dividing \( | G  |  \). Then \( H \trianglelefteq G  \). Note that this powerful idea makes one of the problems fro the exam trivial. That is, recall that the problem was:   
\begin{center}
    \( | G  |  = p^{3} q  \) with \( p > q  \). If \( H \leq G  \) with \( |  H  |  = p^{3}   \) . Then \( | G: H  | = q  \) which is the smallest prime which divides \( | G  |  \), \( H \trianglelefteq G  \).
\end{center}

\subsection{Conjugacy in \( {S}_{n} \)}

Recall that the important idea is not the permutation representation but rather the conjugacy classes which come from \( G  \) acting on itself by conjugation. Let \( \sigma = (1 \ 2 \ 3) \in {S}_{5} \) and let \( \tau = (3  \ 5 ) \in {S}_{5} \). Then
\[  \tau \sigma \tau^{-1} = (3 \ 5)(1 \ 2 \ 3)(3 \ 5) = (1 \ 2 \ 5)    \]
and
\[  \tau \sigma \tau^{-1} = (2 \ 3) (1 \ 2 \ 3)(2 \ 3)^{-1} = (1 \ 3 \ 2) \]
and
\begin{align*}
    \tau \sigma \tau^{-1} &= (1 \ 4 \ 3 \ 2) (1 \ 2 \ 3)(1 \ 4 \ 3 \ 2)^{-1} \\
                          &= (4 \ 1 \ 2)
\end{align*}

\begin{lemma}
    Let \( \sigma = ({a}_{1} \ {a}_{2} \ \dots \ {a}_{k}) \in {S}_{n} \) be a \( k \)-cycle, then for any \( \tau \in {S}_{n} \), we have \[ \tau \sigma \tau^{-1} = (\tau({a}_{1}), \tau({a}_{2}), \tau({a}_{3}), \dots, \tau({a}_{k})).  \]
\end{lemma}
\begin{proof}
Notice that since \( \sigma({a}_{1}) = {a}_{2} \), then \( \tau \sigma \tau^{-1} ( \tau({a}_{1})) = \tau \sigma ({a}_{1}) = \tau({a}_{2})  \) Similarly, \( \sigma({a}_{2}) = {a}_{3} \), and so  
\[  \tau^{-1}(\tau({a}_{2}))  = \tau({a}_{3}).\]
In general, \( \sigma({a}_{i}) = {a}_{i+1} \) and so \( \tau \sigma \tau^{-1} (\tau({a}_{i})) = \tau({a}_{i+1}) \). For \( 1 \leq i < k  \). Finally \( \sigma({a}_{k}) =  {a}_{1} \) and so  
\[  \tau \sigma \tau^{-1} (\tau({a}_{k})) = \tau({a}_{1}). \]
Noting that all the \( \tau({a}_{i}) \) are distinct since \( \tau  \) is injective and all the \( {a}_{i} \)'s are distinct, we have \( \tau \sigma \tau^{-1} \) contains the \( k  \)-cycle 
\[  (\tau({a}_{1}), \tau({a}_{2}), \dots, \tau({a}_{k})). \]
Let \( b \in A = \{  1, 2, \dots, n \}  \). But \( b \notin \{ {a}_{1}, {a}_{2}, \dots, {a}_{k} \}  \), then \( \sigma(b) = b \) and it follows that \( \tau \sigma \tau^{-1} (\tau(b)) = \tau(b) \); that is, we see that \( \tau(b) \) is fixed. Hence,  
\[  \tau \sigma \tau^{-1} = (\tau({a}_{1}), \tau({a}_{2}), \dots \tau({a}_{k})). \]
\end{proof}

\begin{prop}
    Let \( \sigma  \) and \( \tau \) be elements of \( {S}_{n} \) and suppose that \( \sigma  \) has cycle decomposition 
    \[  {\sigma}_{1} \ {\sigma}_{2} \ {\sigma}_{3} \dots \ {\sigma}_{i} = ({a}_{1} \ a_{2} \ {a}_{3} \dots {a}_{{k}_{i}}),  \]
    then \( \tau \sigma \tau^{-1} \) has cycle decomposition
    \[  {\tau}_{1} \ {\tau}_{2} \ \dots {\tau}_{j} \]
    where \( {\tau}_{i} = (\tau({a}_{1}) \ \tau({a}_{2}) \ \dots \tau({a}_{{k}_{i}})) \).
\end{prop}
\begin{proof}
The proof to the above proposition follows from the fact that \( \tau {\sigma}_{1} {\sigma}_{2} \tau^{-1} = (\tau {\sigma}_{1} \tau^{-1}) (\tau {\sigma}_{2} \tau^{-1}) \). 
\end{proof}

\begin{eg}
   Let \( \sigma = (1 \ 2)(3 \ 4 \ 5 )(6 \ 7 \ 8 \ 9) \) and let \( \tau = (1 \ 3 \ 5 \ 7)(2 \ 4 \ 6 \ 8) \). Then 
   \begin{align*}
       \tau \sigma \tau^{-1} &= (3 \ 4)(5 \ 6 \ 7)(8 \ 1 \ 2 \ 9) \\
   \end{align*}
   Note that if we change the cycle types then the resulting will retain the same cycle type.
\end{eg}

\begin{definition}[ ]
   \begin{enumerate}
       \item[(1)] If \( \sigma \in {S}_{n} \) is the product of disjoint cycles of lengths \( {n}_{1}, {n}_{2}, {n}_{3}, \dots, {n}_{r}  \) where \( {n}_{1} \leq {n}_{2} \leq \dots \leq {n}_{r} \) (including \( 1  \)-cycles), the integers \( {n}_{1}, {n}_{2}, \dots, {n}_{r} \) are called the cycle type of \( \sigma \).
       \item[(2)]
    If \( n \in \Z^{+} \), a partition of \( n  \) is any non-decreasing sequence of positive integers whose sum is \( n  \). 
   \end{enumerate} 
\end{definition}

The punchline here is that cycle types have a one-to-one correspondence with partitions of \( n  \) \textbf{(prove it)}.

We have seen that conjugation preserves cycle type; that is, two elements of \( {S}_{n} \) are conjugate implies they have the same cycle type. Consider the converse of this statement; that is, if two elements of \( {S}_{n} \) have the same cycle type, then they are conjugates of each other. Note that  
\[  {\sigma}_{1} = (3 \ 4 \ 7 \ 2)(1 \ 5 \ 8) \]
and
\[  {\sigma}_{2} = (1 \ 2 \ 3 \ 4)(5 \ 6 \ 7). \]
We can find \( \tau \) such that \( \tau {\sigma}_{1} \tau^{-1} = {\sigma}_{2}. \)
Define \( \tau  \) as \( \tau(3) = 1  , \tau(4) = 2  , \tau(7) = 3 , \tau(2) = 4, \tau(1) = 5 , \tau(5) = 6, \tau(6) = 8 \) and \( \tau(8) = 7  \).
Then 
\[  \tau = (1 \ 5 \ 6  \ 8 \ 7 \ 3)(2 \ 4) \]

\begin{prop}[Proposition 11 from the book]
   Two elements of \( {S}_{n} \) are conjugate if and only if they have the same cycle type. The number of cycle types of \( {S}_{n} \) equals the number of partitions of \( n  \).
\end{prop}

Note that from this we see that each cycle type generates some kind of equivalence class. The goal of next lecture is to show that \( {A}_{n} \) is a simple group.

\subsection{Looking at the Conjugacy Class of \( {S}_{5} \)}
Last time we learned that each conjugacy class in \( {S}_{n} \) consists of all elements in \( {S}_{n} \) with the same cycle type. Let's look at \( {S}_{5} \). What are the conjugacy classes of \( {S}_{5} \)? All elements with cycle type. Let's identify them. 

Say, we have \( 1, 1, 1, ,1 ,1  \) corresponds to the representative \( {1}_{{S}_{5}} \). If it is a 5 cycle, then we have \( (1 \ 2 \ 3 \ 4 \ 5) \) as a representative. We can do the same process for all the other cycle types in \( {S}_{5} \).    

What are the number of elements in the conjugacy class of \( 1  \)? There are \( 1  \) one of them, namely, the identity \( {1}_{{S}_{5}} \). For \( 1 , 1 , 1, 2  \), the number of elements in ht conjugacly classes is \( 1 0 \) because   
\[  \frac{ 5 \cdot 4  }{ 2 }  = 10 \]
and for \( 1,2, 3 \), the number of elements in the conjugacy classes is  
\[  \frac{ 5 \cdot 4 \cdot 3  }{ 3 } = 2 0.   \]
We can do a similar process for the rest of the cycle types in \( {S}_{5} \). 

We can use the number of elements in a conjugacy class to find the centralizer of an element in \( {S}_{5} \). Let \( \sigma = (1 \ \ 2 \ 3) \). Denote the conjugacy class of \( \sigma \) as \( {K}_{\sigma} \) which has \[ | {K}_{\sigma} |  = 20 = | G : {C}_{{S}_{5}}(\sigma) | = \frac{ | {S}_{5} | }{ | {C}_{{S}_{5}}(\sigma) |  }  = \frac{ 5 !  }{ | {C}_{{S}_{5}} (\sigma)|  }.    \] 
This tells us that 
\[ \frac{ 5!  }{  20  }  = | {C}_{{S}_{5}}(\sigma) | \implies 3! = | {C}_{{S}_{5}} (\sigma) |.     \]
Thus, we see that \( | {C}_{{S}_{5}}(\sigma) |  = 6  \).

Note that we know \( \langle  \sigma \rangle \leq {C}_{{S}_{5}}(\sigma) \). We know that
\[  {C}_{{S}_{5}}(\sigma) = \{  1, \sigma , \sigma^{2} \}  \]
We also know that disjoint cycles commute so \( (4 \ 5) \) commutes with \( \sigma \). Then 
\[  {C}_{{S}_{5}}(\sigma) = \{  1 , \sigma, \sigma^{2} , (4 \ 5) , (4 \ 5) \sigma, (4 \ 5) \sigma^{2}  \}.  \]
Find \( {C}_{{S}_{5}}((1 \ 2 \ 3 \ 4 \ 5)) \) in \( {S}_{5} \).

% fill this out later


More generally, for \( \sigma  \) an \( m  \)-cycle in \( {S}_{n} \), we have \( {K}_{\sigma}  \) has
\[  \frac{ n(n-1) \cdots (n - (m - 1))  }{ m  } = \frac{ n!  }{ m (n - m)! }   \]
where \( m  \) is the number of representations of the \( m  \)-cycle. So, we have that 
\[  \frac{ | {S}_{n} |   }{ | {C}_{{S}_{n}}(\sigma) |  }  = \frac{ n!  }{ m (n-m)! }. \]
This tells us further that 
\[  | {C}_{{S}_{n}}(\sigma) |  = m (n -m )! \]
Note that there are \( m \) elements in \( \langle  \sigma \rangle \) and there are \( (n-m)! \) elements of \( {S}_{n} \) which fix \( 1, 2, 3, \dots, m  \). We can look at this in the following way: for \( n > m  \), the first spot past \( m  \), you will have \(  n - m \) choices to fix \( 1, 2, 3, \dots, m  \). For \(  m + 1  \), you will have \(  (n -m ) -  1  \) choices to fix the elements of the \( m  \) cycle. Generalize this until you get to \( n - 1  \), then we have we will only have \( 2 \) choices and so by the time we get to \( n  \), we will only have one choice left to make.


So, \( {C}_{{S}_{n}}(\sigma) \) consists of \( \sigma^{i} \tau  \) where \( j = 0 , \dots, m - 1  \) and \( \tau  \) is a permutation that is disjoint from \( \sigma \). Since there are \( m (n-m)! \)  elements of this form they make up all elements of \( {C}_{{S}_{n}}(\sigma) \).

\subsection{Proving that the Alternating Group is Simple}

We will prove hat \( {A}_{5} \) is simple, bu we need a general fact. Here it is:
\begin{prop}[A normal subgroup is the union of its conjugacy classes]
    For any group \( G  \) and \( H \trianglelefteq G  \), we have 
    \[  H = \bigcup_{ h \in H  }^{   }  {K}_{h}. \]
    That is \( H  \) is the union of a collection of conjugacy classes.
\end{prop}
\begin{proof}
Let \( h \in H  \). Then \( h  \) is in exactly one conjugacy class, namely, \({K}_{h} \). Hence,   
\[  h \in \bigcup_{ h \in H  }^{  }  {K}_{k }. \]
Thus, 
\[  H \subseteq  \bigcup_{ h \in H  }^{  }  {K}_{h}. \]
Let \( k \in {K}_{h} \) for some \( h \in H  \). Then \( k = g h g^{-1} \) for some \( g \in G  \). Since \( H  \) is normal in \( G  \), \( k \in H  \). Thus 
\[  H = \bigcup_{  h \in H  }^{   }  {K}_{h}. \]
\end{proof}

The consequence of this fact is: If \( |  H  | = n  < \infty   \), then 
\[  | H  |  = |{K}_{1}| + |{K}_{2}| + \cdots + |{K}_{j}| \]
where the \( {K}_{i} \) as a collection of conjugacy classes of \( G  \). Now let's prove that \( {A}_{5} \) is simple. We will do this by finding the order of all conjugacy classes and show they cannot add up to any divisor of the \( | {A}_{5} | = \frac{ 5 !  }{  2  }  = 60   \) except \( | {A}_{5} |   \) and \( 1  \).     

Let's look at the representatives of \( {S}_{5}  \); namely, \( 1  \), \( (1 \ 2)  \), \( (1 \ 2 \ 3 \ 4) \), \( (1 \ 2 \ 3 \ 4 \ 5) \), \( (1 \ 2 ) (3 \ 4) \). Note that \( 1 \ 2  \), \( (1 \ 2 \ 3) \) and \( (1 \ 2 ) ( 3 \ 4 \ 5) \) are odd. Let's start finding the size of the conjugacy classes of \(  1 , (1 \ 2 \ 3) \), and \( ( 1 \ 2 \ 3 \ 4 \ 5) \) in \( {A}_{5} \). Note that  
\[  {K}_{1} = \{ 1  \}.  \]
\[ | {K}_{(1 \ 2 \ 3 )}   |  = \frac{ | {A}_{5} |    }{  | {C}_{{A}_{5}}((1 \ 2 \ 3)) |  }  = \frac{ \frac{ 5!  }{  2  }  }{ 3 }  \]
where 
\[  {C}_{{A}_{5}} = \{ 1 , (1 \ 2 \ 3), (1 \  2 \ 3 )^{2},  \}  \]
but note that \( (1 \ 2 \ 3)^{2} ( 4 \ 5), (4 \ 5), (1 \ 2 \ 3 )(4 \ 5)  \) are not in \( {A}_{5} \) because they are odd.
Hence, we see that \( | {K}_{(1 \ 2 \ 3)} |  = 2 0 \). All \( 3  \) cycles in \( {A}_{5} \) are conjugate of each other. Now, let's look at \[ | {K}_{(1 \ 2 \ 3 \ 4 \ 5)  } |  = \frac{ | {A}_{5} |   }{ | {C}_{{A}_{5}}((1 \ 2 \ 3 \ 4 \ 5)) |  }.  \]
Since all powers of a \( 5  \)-cycle are even, we have 
\[  | {C}_{{A}_{5}}(1 \ 2 \ 3 \ 4 \ 5) |  = 5.  \]
This gives us 
\[  | {K}_{(1 \ 2 \ 3 \ 4 \ 5)} | = \frac{ \frac{ 5!  }{  2  }  }{  5 }.  \]
Notice there are \( 2 4  \) \( 5 \)-cycles in \( {A}_{5} \). Thus, there must be a \( 5- \)cycle \( \sigma \) such that \( \sigma  \) is not conjugate to \( (1 \ 2 \ 3 \ 4  \ 5) \). Another example is \( ( 1 \ 2 \ 5 \ 3 \ 4) \). The exact same calculation gives that \( | {K}_{\sigma} | = 12  \). Lastly, let's figure the size of \( {K}_{(1 \ 2 )( 3 \ 4)} \). We know that \( |{K}_{(1 \ 2)(3 \ 4)} | = \frac{ | {A}_{5} |  }{ | {C}_{{A}_{5}} ((1 \ 2 )(3 \ 4)) |  }  \). Notice that \( {C}_{{A}_{5}}((1 \ 2 )(3 \ 4)) \) is a \textbf{group} and  
\[  {C}_{{A}_{5}}((1 \ 2)(3 \ 4)) \subseteq  {C}_{{S}_{5}}((1 \ 2)(3 \ 4)). \]
Thus, 
\[  {C}_{{A}_{5}}((1 \ 2)(3 \ 4)) \leq {C}_{{S}_{5}}((1 \ 2)(3 \ 4)) \]
which implies its order divides \( 8  \).

This implies that all cycle types \( 1 , 2, 2  \) in \( {S}_{5} \) are conjugate as \( | {K}_{(1 \ 2)(3 \ 4)} = 15 |  \). So a normal subgroup of \( {A}_{5} \) must have order \( 1  \). The sum of a collection \( 1, 20, 12, 12, 15  \) but none of these sums divide except \( 1  \) and \( 60  \). Thus, \( {A}_{5} \) is simple.


Can you explain the following argument to cayley's theorem, Thus, \(k \mid \frac{ p!  }{  p  }  = (p-1)! \). But all prime divisors of \( (p-1)! \) are less than \( p  \) and by the minimality of \( p  \), every prime divisor of \(  k  \) is greater than or equal to \( p  \). This forces \( k = 1  \), so \( H = K \trianglelefteq G  \). Here is is assumed that \( p  \) is the smallest prime dividing \(  | G  |  \).
